% p.41 — الأفعال المعتلة (تابع): الأعمدة (25)–(30)
% Layout: six columns, each headed by a circled number, the verb type, then the bāb
% (e.g. حَسِبَ). Each column has two blocks of ten rows (ماضي، مضارع، مصدر، اسم فاعل،
% ماضي مجهول، مضارع مجهول، مصدر، اسم مفعول، أمر، نهي), separated by a rule: the upper
% block gives the uncontracted (أصل) forms, the lower block the forms after إعلال.
\mspage{41}
\begin{sarf}{|c|c|c|c|c|c|}
\thead{(25) مثال يائي} & \thead{(26) مثال يائي} & \thead{(27) مثال ومهموز} & \thead{(28) اجوف واوي} & \thead{(29) اجوف واوي} & \thead{(30) اجوف واوي}\\\hline
حَسِبَ & كَرُمَ & ضَرَبَ & نَصَرَ & فَتَحَ & سَمِعَ\\\hline
يَقِظَ & يَمُنَ & وَأَرَ & قَوَلَ & بَوَءَ & خَوِفَ\\
يَيْقَظُ & يَيْمُنُ & يَوْئِرُ & يَقْوُلُ & يَبْوَءُ & يَخْوَفُ\\
يَقَظَةً & يُمْنًا & وَأْرًا & قَوْلًا & بَوْءًا & خَوْفًا\\
يَاقِظٌ & يَامِنٌ & وَائِرٌ & قَاوِلٌ & بَاوِءٌ & خَاوِفٌ\\
\x & \x & وُئِرَ & قُوِلَ & بُوِءَ & خُوِفَ\\
\x & \x & يُوْئَرُ & يُقْوَلُ & يُبْوَءُ & يُخْوَفُ\\
\x & \x & وَأْرًا & قَوْلًا & بَوْءًا & خَوْفًا\\
\x & \x & مَوْئُوْرٌ & مَقْوُوْلٌ & مَبْوُوْءٌ & مَخْوُوْفٌ\\
اِيْقَظْ & اُيْمُنْ & اِوْئِرْ & اُقْوُلْ & اِبْوَءْ & اِخْوَفْ\\
لَا تَيْقَظْ & لَا تَيْمُنْ & لَا تَوْئِرْ & لَا تَقْوُلْ & لَا تَبْوَءْ & لَا تَخْوَفْ\\\hline
يَقِظَ & يَمُنَ & وَأَرَ & قَالَ & بَاءَ & خَافَ\\
يَيْقَظُ & يَيْمُنُ & يَئِرُ & يَقُوْلُ & يَبَاءُ & يَخَافُ\\
يَقَظَةً & يُمْنًا & وَأْرًا & قَوْلًا & بَوْءًا & خَوْفًا\\
يَقْظَانُ & يَمِيْنٌ & وَائِرٌ & قَائِلٌ & \unclear{بَائِءٌ} & خَائِفٌ\\
\x & \x & وُئِرَ & قِيْلَ & بِيْءَ & خِيْفَ\\
\x & \x & يُوْئَرُ & يُقَالُ & يُبَاءُ & يُخَافُ\\
\x & \x & وَأْرًا & قَوْلًا & بَوْءًا & خَوْفًا\\
\x & \x & مَوْئُوْرٌ & مَقُوْلٌ & مَبُوْءٌ & مَخُوْفٌ\\
اِيْقَظْ & \unclear{اُوْمُنْ} & \unclear{اِيْئِرْ ءِرْ} & قُلْ & بَءْ & خَفْ\\
لَا تَيْقَظْ & لَا تَيْمُنْ & لَا تَئِرْ & لَا تَقُلْ & لَا تَبَءْ & لَا تَخَفْ\\
\end{sarf}
